% =====================================================
% 题型：立体几何多结论选择题 (q_solid_sphere_properties.tex)
% =====================================================
% =====================================================
% 难度：★★★ (难题)
% =====================================================
\newcommand{\qSolidSpherePropertiesStem}{如图，\(AB\) 是 \(\odot O\) 的直径，\(PA \perp\) 平面 \(\odot O\) 所在的平面，\(C\) 是圆周上不同于 \(A,B\) 的任意一点，则 \hfill （\quad）}

\newcommand{\qSolidSpherePropertiesOptions}{%
  \mcoptions[1]{
    若 $D$ 为 $BC$ 的中点，则 $OD \parallel$ 平面 $PAC$
  }{
    平面 $PAC \perp$ 平面 $PBC$
  }{
    点 $A$ 到平面 $PBC$ 的距离小于 $\dfrac{1}{2} PC$
  }{
    $PB$ 为三棱锥 $P-ABC$ 外接球的直径
  }%
}

\newcommand{\qSolidSpherePropertiesFigure}[1][1.2]{%
  \begin{tikzpicture}[scale=#1, >=stealth]
    \coordinate (O) at (1.2, 0);
    \coordinate (A) at (0, 0);
    \coordinate (B) at (2.4, 0);
    \coordinate (P) at (0, 1.8);
    \coordinate (C) at (0.7, 0.25);
    \coordinate (O1) at (1.2, 1.8);
    \coordinate (Ctop) at (0.7, 2.05);
    
    \draw[thick] (2.4,0) arc (0:-180:1.2 and 0.35);
    \draw[dashed] (2.4,0) arc (0:180:1.2 and 0.35);
    
    \draw[dashed] (1.2,1.8) ellipse (1.2 and 0.35);
    \draw[dashed] (2.4,0) -- (2.4,1.8);
    \draw[dashed] (O) -- (O1);
    \draw[dashed] (C) -- (Ctop);
    
    \draw[thick] (P) -- (A);
    \draw[thick] (P) -- (B);
    \draw[thick, dashed] (P) -- (C);
    \draw[thick, dashed] (A) -- (C);
    \draw[thick, dashed] (C) -- (B);
    \draw[thick] (A) -- (B);
    
    \node[left] at (A) {\small $A$};
    \node[right] at (B) {\small $B$};
    \node[above left] at (P) {\small $P$};
    \node[below] at (O) {\small $O$};
    \node[below right] at (C) {\small $C$};
  \end{tikzpicture}%
}

\newcommand{\qSolidSpherePropertiesAnswer}{A, B, D}

\newcommand{\qSolidSpherePropertiesSolution}{%
  因为 \(PA\perp\) 平面 \(\odot O\)，且 \(AB\subset\) 平面 \(\odot O\)，所以 \(PA\perp AB\)，即 \(\angle PAB=90^\circ\)。\\
  又因为 \(AB\) 是圆 \(O\) 的直径，\(C\) 在圆周上，由圆周角定理可知 \(AC\perp BC\)。\\
  由于 \(PA\perp\) 平面 \(\odot O\)，且 \(BC\subset\) 平面 \(\odot O\)，得 \(PA\perp BC\)。\\
  因为 \(PA \cap AC = A\)，且 \(PA, AC \subset\) 平面 \(PAC\)，所以 \(BC\perp\) 平面 \(PAC\)。\\
  因为 \(BC \subset\) 平面 \(PBC\)，所以平面 \(PAC \perp\) 平面 \(PBC\)（B 正确）。\\
  因为 \(BC\perp\) 平面 \(PAC \implies BC\perp PC \implies \angle PCB=90^\circ\)。\\
  于是 \(\triangle PAB\) 与 \(\triangle PCB\) 均为直角三角形，且共斜边 \(PB\)。\\
  取 \(PB\) 中点 \(M\)，由直角三角形斜边中点到各顶点距离相等，有 \(MA = MB = MC = MP = \dfrac{1}{2}PB\)。\\
  故 \(PB\) 为三棱锥 \(P-ABC\) 外接球的直径，中点 \(M\) 为球心（D 正确）。\\
  若 \(D\) 是 \(BC\) 的中点，且 \(O\) 是 \(AB\) 的中点，则由三角形中位线定理得 \(OD \parallel AC\)。\\
  因为 \(OD \not\subset\) 平面 \(PAC\)，且 \(AC \subset\) 平面 \(PAC\)，所以 \(OD \parallel\) 平面 \(PAC\)（A 正确）。\\
  对于 C，因为平面 \(PAC \perp\) 平面 \(PBC\)，交线为 \(PC\)。过点 \(A\) 作 \(AH \perp PC\) 于 \(H\)，则 \(AH \perp\) 平面 \(PBC\)，即 \(AH\) 为点 \(A\) 到平面 \(PBC\) 的距离。
  在 Rt\(\triangle PAC\) 中，点 \(A\) 到 \(PC\) 的距离为 \(AH = \dfrac{PA \cdot AC}{PC}\)。
  比较 \(AH\) 与 \(\dfrac{1}{2}PC\)：
  因为 \(PA^2 + AC^2 \ge 2PA \cdot AC\)，且 \(PA^2 + AC^2 = PC^2\)，
  所以 \(PC^2 \ge 2PA \cdot AC \implies \dfrac{PA \cdot AC}{PC} \le \dfrac{1}{2}PC\)，即 \(AH \le \dfrac{1}{2}PC\)，
  其中当且仅当 \(PA = AC\) 时等号成立。
  因此，当 \(PA = AC\) 时，点 \(A\) 到平面 \(PBC\) 的距离恰好等于 \(\dfrac{1}{2}PC\)，
  “点 \(A\) 到平面 \(PBC\) 的距离小于 \(\dfrac{1}{2} PC\)”不能恒成立，故 C 错误。
  综上，结论正确的有 ABD。
  故选 ABD。%
}
